\documentclass[11pt, a4paper, twoside]{article}

% -------------------------------------------------------------------
% PACKAGES — UQU Format
% -------------------------------------------------------------------
\usepackage[utf8]{inputenc}
\usepackage{graphicx}
\usepackage{setspace}
\usepackage{geometry}
\usepackage{mathptmx}
\usepackage{titlesec}
\usepackage{fancyhdr}
\usepackage{tocloft}
\usepackage{float}
\usepackage{url}
\usepackage[colorlinks=true, urlcolor=blue, linkcolor=black, citecolor=black]{hyperref}

% -------------------------------------------------------------------
% MARGINS & SPACING
% -------------------------------------------------------------------
\geometry{top=1.5in, bottom=1.0in, left=1.25in, right=1.0in}
\onehalfspacing
\setlength{\parskip}{6pt}

% -------------------------------------------------------------------
% ADDITIONAL PACKAGES
% -------------------------------------------------------------------
\usepackage{amsmath,amssymb}
\usepackage{booktabs}
\usepackage{longtable}
\usepackage{array}
\usepackage{xcolor}
\usepackage{enumitem}
\usepackage{caption}
\usepackage{subcaption}
\usepackage{multirow}
\usepackage[numbers,sort&compress]{natbib}

% -------------------------------------------------------------------
% HEADINGS (Article class)
% -------------------------------------------------------------------
\titleformat{\section}
  {\normalfont\fontsize{16}{19}\selectfont\bfseries\centering}
  {\thesection}{1em}{\MakeUppercase}

\titleformat{\subsection}
  {\normalfont\fontsize{14}{17}\selectfont\bfseries}
  {\thesubsection}{1em}{}

\titleformat{\subsubsection}
  {\normalfont\fontsize{13}{16}\selectfont\bfseries\itshape}
  {\thesubsubsection}{1em}{}

% -------------------------------------------------------------------
% HEADERS & FOOTERS
% -------------------------------------------------------------------
\pagestyle{fancy}
\fancyhf{}
\fancyhead[LO]{THAMAN}
\fancyhead[RE]{\leftmark}
\fancyfoot[R]{\thepage}
\renewcommand{\headrulewidth}{0.4pt}
\renewcommand{\footrulewidth}{0pt}
\fancypagestyle{plain}{\fancyhf{}\renewcommand{\headrulewidth}{0pt}}
\setlength{\headheight}{14pt}

% -------------------------------------------------------------------
% CUSTOM STYLES & COMMANDS
% -------------------------------------------------------------------
\newcommand{\textarabic}[1]{#1}

\begin{document}

% ── Title page ────────────────────────────────────────────────────────────────
\begin{titlepage}
    \fontfamily{ptm}\selectfont
    \noindent
    \begin{minipage}[t]{0.6\textwidth}
        \raggedright \itshape
        BSc Project\\
        Dept.\ of Computer Science \& Artificial Intelligence\\
        Project ID: CSAI-472-P2-M20\\
        2026
    \end{minipage}%
    \begin{minipage}[t]{0.4\textwidth}
        \raggedleft
        \includegraphics[width=3cm]{Umm_Al-Qura_University_logo.png}
    \end{minipage}

    \vspace{2cm}
    \begin{center}
        {\fontsize{26}{31}\selectfont \textbf{THAMAN}}\\[0.8cm]
        {\fontsize{14}{16}\selectfont \textbf{A Dual-City Automated Valuation Model\\
        for New York City and Riyadh Using\\
        Ensemble Machine Learning and Quality-of-Life Spatial Indicators}}\\[1cm]
        {\fontsize{14}{16}\selectfont \textbf{By:\\[0.3cm]
        Turki AL-murahhem\\}}
    \end{center}

    \vfill
    \noindent
    \begin{flushleft}
        Dept.\ of Computer Science and Artificial Intelligence\\
        Faculty of Computer and Information Systems\\
        Umm Al-Qura University, KSA\\[0.3cm]
        {\small Submitted in partial fulfilment of the requirements for the degree of
        Bachelor of Science in Computer Science.}
    \end{flushleft}
\end{titlepage}

\newpage\thispagestyle{empty}\mbox{}\newpage

% ── Abstract ──────────────────────────────────────────────────────────────────
\begin{abstract}
\noindent
This paper presents \textbf{THAMAN}, an Automated Valuation Model (AVM) for
estimating residential property prices across New York City's five boroughs.
The system integrates structural property attributes with spatial
Quality-of-Life (QoL) indicators---including proximity to transit, crime rates,
income demographics, points-of-interest density, building health signals, and
MTA station quality---into a four-model stacked ensemble comprising two XGBoost
models, LightGBM, and CatBoost, blended by a Ridge meta-learner.

Trained on 185,092 NYC property sales from 2022 to 2026, the final model
(Stack v22) achieves $R^2 = 0.6495$ and Median Absolute Percentage Error
(MedAPE) of 20.32\% on a time-based holdout set of 27,763 unseen sales across
134 features, including a novel NTA price trend slope feature capturing
neighbourhood-level price appreciation momentum via per-NTA OLS regression,
and ten new building health and mobility indicators sourced from NYC HPD
violation records, DOB permit filings, granular 311 complaint density (rodent
and heat), and MTA station quality attributes. The system implements quality controls aligned with the 2025 Interagency
AVM Quality Control Rule (Dodd-Frank \S 1473, effective October 2025),
including segment-adaptive confidence intervals, a 0--100 confidence score
with letter grade, comparable-count hit-rate signalling, and four quality
control flags (\texttt{SPARSE\_MARKET}, \texttt{LUXURY\_SEGMENT},
\texttt{HIGH\_UNCERTAINTY}, \texttt{METRO\_CORE}).

A bilingual (English/Arabic) web application provides interactive map-based
property valuation with SHAP-driven feature explanations including waterfall
visualisations, natural-language tooltips, a segment-adaptive confidence band,
and real-time neighbourhood-level encoding resolved from a 185K-sale spatial
index. The complete system is deployed on Hugging Face Spaces and open-sourced
on GitHub.

The system is extended to Riyadh, Saudi Arabia as a parallel AVM trained on
6,910 district-level real estate transactions (2018--2025). Across 140 features
covering metro and bus transit, commercial density, air quality
(NO$_2$, SO$_2$, PM$_{10}$, O$_3$), QoL proximity to mosques, malls,
schools, hospitals, parks and entertainment venues, macroeconomic indices (REI
quarterly index, Saudi average salary), district price aggregates, and rental
market signals (SA\_Aqar), the Riyadh stack achieves OOF $R^2 = 0.9343$ /
MedAPE~$= 8.28\%$ and holdout $R^2 = 0.8003$ / MedAPE~$= 15.56\%$ on a 2025
Q1--Q3 out-of-sample test set of 1,727 district-quarter observations.

\medskip\noindent
\textbf{Keywords:} automated valuation model, ensemble learning, XGBoost,
LightGBM, CatBoost, stacking, target encoding, spatial features, SHAP, New York
City, Riyadh, Saudi Arabia, dual-city AVM, real estate, quality of life
indicators, adaptive confidence, NTA encoding, district-level prediction.
\end{abstract}

\tableofcontents
\newpage

% ─────────────────────────────────────────────────────────────────────────────
\section{Introduction}
\label{sec:intro}

\subsection{Motivation}
Property valuation is a central task in real estate finance, mortgage
underwriting, and urban policy. Traditional appraisal is slow (days to weeks),
expensive (\$300--\$500 per appraisal), and subject to appraiser bias.
Automated Valuation Models (AVMs) use statistical and machine learning methods
to produce instant price estimates at near-zero marginal cost.

The deployment of AVMs in highly heterogeneous real estate markets presents a
formidable computational challenge at the intersection of data science, urban
economics, and municipal policy. Urban environments characterised by dense
stratification, hyper-local zoning regulations, and vast disparities in
structural typologies require valuation frameworks that transcend simple linear
hedonic pricing or naive comparative market analyses.

Existing commercial AVMs (Zillow Zestimate, Redfin Estimate, CoreLogic) are
proprietary and opaque. Empirical evaluations indicate that while they modestly
outperform naive list prices, they remain highly vulnerable to geographic
heterogeneity and tail risks. Furthermore, algorithmic bias presents a systemic
risk: analyses demonstrate that AVMs can systematically yield greater valuation
errors for minority homeowners, producing errors significantly higher in
majority-Black neighbourhoods even after controlling for structural
characteristics.

This project addresses four gaps:
\begin{enumerate}
  \item Most NYC-focused models use small datasets (5K--50K sales). This work
        uses 185K sales across five boroughs and four years.
  \item QoL spatial indicators are rarely integrated end-to-end from raw open
        data into a production-ready inference pipeline.
  \item AVM confidence communication is often binary. This work implements
        segment-adaptive confidence with a formal QC block aligned with 2026
        AVM industry standards.
  \item Most models ignore neighbourhood price momentum. This work introduces
        an NTA-level log-linear price trend slope as an explicit feature.
\end{enumerate}

\subsection{Research Questions}
\begin{description}[style=nextline]
  \item[RQ1] Can a stacked ensemble of gradient boosted trees outperform a
             single XGBoost baseline on NYC residential property price
             prediction?
  \item[RQ2] Do neighbourhood-level spatial QoL features (NTA encoding,
             transit proximity, crime rate, income) improve prediction accuracy
             over structural attributes alone?
  \item[RQ3] How should confidence intervals be communicated to reflect the
             heterogeneous predictability across boroughs and price tiers?
  \item[RQ4] Does incorporating neighbourhood price appreciation momentum
             (NTA price trend slope) further improve model accuracy?
  \item[RQ5] Can the same stacking architecture generalise to a data-scarce
             emerging market (Riyadh) where observations are district-level
             aggregates rather than individual transactions, and how does the
             OOF-to-holdout gap behave when the holdout represents entirely
             unseen calendar quarters?
\end{description}

\subsection{Contributions}
\begin{itemize}
  \item A full end-to-end AVM pipeline from raw NYC open data to deployed web
        application.
  \item A 134-feature dataset combining structural, spatial, temporal,
        neighbourhood momentum, building health, and transit quality signals.
  \item Stack v22: a 4-model ensemble achieving $R^2 = 0.6495$ / MedAPE
        $= 20.32\%$ on 27,763 holdout sales, with 10 new building health
        and mobility features (HPD violations, DOB permits, 311 rodent/heat
        density, MTA station quality attributes).
  \item Segment-adaptive confidence intervals per borough $\times$ price tier.
  \item A formal AVM QC block with four risk flags.
  \item A bilingual (English/Arabic) interactive valuation interface with SHAP
        waterfall plots, natural-language explanations, and 11-layer NTA
        choropleth.
  \item Riyadh Stack v11: a dual-city extension to Saudi Arabia using 6,910
        district-quarter observations, 140 features, and the identical stacking
        architecture. Achieves holdout $R^2 = 0.8003$ / MedAPE $= 15.56\%$ on
        a fully out-of-sample 2025 Q1--Q3 test set (1,727 rows).
\end{itemize}

% ─────────────────────────────────────────────────────────────────────────────
\section{Related Work}
\label{sec:related}

\subsection{The Evolution and Limitations of Traditional AVMs}
Automated valuation algorithms are software programs that market analysts,
lenders, and institutions use to produce market value estimates based on
location, market conditions, and real estate characteristics. From their
inception, they have been utilised by government-sponsored enterprises (GSEs)
and financial institutions for underwriting lower loan-to-value mortgages and
portfolio-wide collateral risk assessments.

The baseline performance of contemporary commercial AVMs in dense urban cores
typically yields a MedAPE between 17.5\% and 19.8\%, with models frequently
overstating administrative assessments by 16--18\%. These models routinely
struggle to maintain predictive parity across diverse property classes, often
generating wider error distributions in heterogeneous zip codes or historically
marginalised neighbourhoods.

\subsection{Hedonic Price Models}
The hedonic pricing framework \cite{rosen1974} decomposes property price into
implicit prices of constituent attributes. Early models used OLS regression on
structural features: floor area, number of rooms, age, and location dummies.
These typically achieve $R^2$ of 0.5--0.65 on urban datasets
\cite{sirmans2005}.

\subsection{Machine Learning for Property Valuation}
Gradient boosted trees have become the dominant method in AVM literature.
XGBoost \cite{chen2016}, LightGBM \cite{ke2017}, and CatBoost
\cite{prokhorenkova2018} consistently outperform OLS and neural networks on
tabular real estate data \cite{phan2018,truong2020}.

Ensemble stacking---training a meta-learner on out-of-fold predictions of
diverse base models---further improves accuracy by reducing correlated errors
\cite{wolpert1992,breiman1996}.

\subsection{Spatial and QoL Features}
Proximity to transit is a well-documented price determinant. A New York MTA
study found subway station proximity adds \$0--\$8K per 100\,m reduction in
walking distance \cite{hess2007}. Park proximity \cite{lutzenhiser2001}, school
quality \cite{black1999}, and crime rates \cite{ihlanfeldt2010} have all been
shown to capitalise into residential prices.

\subsection{AVM Confidence and Quality}
The 2025 Interagency Final Rule on AVM Quality Control Standards
\cite{avm_qc_rule2025} (Dodd-Frank \S 1473, effective 1 October 2025)
mandates that lending institutions deploying AVMs for mortgage decisions
implement policies addressing credibility, data integrity, conflict-of-interest
avoidance, sample testing, and nondiscrimination. Industry guidance further
recommends that every AVM output include: (a) a confidence score,
(b) a comparable count (hit rate), and (c) explicit quality flags for
data-sparse or atypical market segments.

% ─────────────────────────────────────────────────────────────────────────────
\section{Data}
\label{sec:data}

\subsection{Primary Dataset: NYC Citywide Rolling Calendar Sales}
\begin{description}
  \item[Source] NYC Open Data --- Citywide Property Sales
  \item[Period] January 2022 -- January 2026
  \item[Records] 185,092 residential sales after cleaning
\end{description}

Cleaning steps included:
\begin{itemize}
  \item Removed sales below \$10,000 (non-arms-length transfers)
  \item Removed sales above \$1,000,000,000 (bulk portfolio transfers)
  \item Removed records missing sale\_date, latitude, longitude, or sale\_price
  \item Removed non-residential building classes
\end{itemize}

\begin{table}[H]
  \centering
  \caption{NYC Sales Distribution by Borough}
  \label{tab:borough_dist}
  \begin{tabular}{lrr}
    \toprule
    Borough & Sales & Share \\
    \midrule
    Manhattan     & 46,621 & 25.2\% \\
    Bronx         & 15,158 &  8.2\% \\
    Brooklyn      & 48,105 & 26.0\% \\
    Queens        & 57,070 & 30.8\% \\
    Staten Island & 18,138 &  9.8\% \\
    \midrule
    \textbf{Total} & \textbf{185,092} & \textbf{100\%} \\
    \bottomrule
  \end{tabular}
\end{table}

Price statistics: median \$820,000; mean \$1,778,977; min \$10,040; max
\$1,080,000,000.

\subsection{PLUTO (Primary Land Use Tax Lot Output)}
PLUTO 25v4 (NYC Department of City Planning) is an exhaustive compilation of
land use, geographic, and building-level data at the tax lot level. Merging on
BBL (Borough-Block-Lot) identifier achieves 98.5\% coverage.

The \texttt{assesstot} field was used to impute \texttt{prior\_sale\_price} via
the assesstot ratio method:
\begin{equation}
  \hat{P}_{\text{prior}} = P_{\text{current}} \times
    \frac{\mathit{assesstot}_{\text{prior}}}{\mathit{assesstot}_{\text{current}}}
  \label{eq:assesstot}
\end{equation}
This achieved 98.0\% coverage on \texttt{prior\_sale\_price}.

\subsection{The Fractional Assessment Labyrinth}
Unlike many jurisdictions that tax properties at full market value, NYC's real
property tax system applies assessment ratios that vary dramatically by tax
class, originating from the 1975 \textit{Hellerstein v. Town of Islip} decision:
\begin{description}
  \item[Class 1] (1--3 family homes): target ratio 6\% of market value
  \item[Class 2] (larger residential): target ratio 45\% of market value
  \item[Class 4] (commercial): target ratio 45\% of market value
\end{description}
State law caps annual assessment growth for Class 1 at 6\%/year (20\% over 5
years) and Class 2 at 8\%/year (30\% over 5 years). By FY 2024 the median
effective assessment ratio for Class 1 properties had degraded to 3.9\%---far
below the 6\% statutory target. The THAMAN imputation calibrates using the
actual observed sale-price-to-assesstot ratio distribution within each borough.

\subsection{Spatial Datasets}

\begin{table}[H]
  \centering
  \caption{Spatial data sources used in feature engineering}
  \label{tab:spatial_sources}
  \begin{tabular}{llr}
    \toprule
    Dataset & Source & Records \\
    \midrule
    Subway stations         & MTA / NYC Open Data     &    472 \\
    Bus stops               & MTA GTFS                & 16,904 \\
    Parks properties        & NYC Parks Dept          &  2,058 \\
    Public schools (HS)     & NYC DOE                 &    427 \\
    Public schools (Elem.)  & NYC DOE                 &    423 \\
    Hospital facilities     & NYC DOHMH               &    180 \\
    Airbnb listings         & Inside Airbnb           & 36,261 \\
    Bike lanes (vertices)   & NYC DOT                 & 44,934 \\
    Waterfront/coastline    & NYC Open Data           & 27,116 \\
    POI (cafés, rest., gyms)& OpenStreetMap           & 60,000+ \\
    MTA station quality     & MTA / NYC Open Data     &    496 \\
    HPD violations by ZIP   & NYC DOHMH Socrata       &    562 groups \\
    DOB permits by ZIP      & NYC Buildings Socrata   &    227 groups \\
    311 rodent/heat         & NYC Open Data           & 162,000+ \\
    \bottomrule
  \end{tabular}
\end{table}

\subsection{Neighbourhood Tabulation Areas (NTA)}
NTAs represent a far more sophisticated spatial taxonomy than ZIP codes. Each
NTA contains a minimum population threshold (15,000--40,000 residents),
mitigating the high sampling error of sparsely populated Census tracts.
Non-residential zones (parks, airports, cemeteries) are explicitly separated.
The THAMAN pipeline uses 212 distinct NTA codes.

\subsection{Macroeconomic Data}
The 30-year fixed mortgage rate was sourced from the Freddie Mac PMMS weekly
survey, matched to each sale by \texttt{sale\_date}. Range in dataset:
3.45\% -- 7.62\%.

% ─────────────────────────────────────────────────────────────────────────────
\section{Feature Engineering}
\label{sec:features}

\subsection{Structural Features}
\begin{table}[H]
  \centering
  \caption{Structural features}
  \label{tab:structural}
  \begin{tabular}{ll}
    \toprule
    Feature & Description \\
    \midrule
    \texttt{gross\_square\_feet} & Total building floor area (sq ft) \\
    \texttt{land\_square\_feet}  & Lot area (sq ft) \\
    \texttt{building\_age}       & Current year $-$ year\_built \\
    \texttt{numfloors}           & Number of floors \\
    \texttt{residential\_units}  & Number of residential units \\
    \texttt{log\_land\_sqft}     & $\log(1+\texttt{land\_sqft})$ \\
    \texttt{lot\_coverage}       & \texttt{gross\_sqft} / \texttt{land\_sqft} \\
    \texttt{bldg\_vol\_proxy}    & $\log(1+\texttt{gross\_sqft}\times\texttt{numfloors})$ \\
    \texttt{prior\_price\_psf}   & prior\_sale\_price / gross\_sqft \\
    \texttt{sqft\_per\_floor}    & gross\_sqft / numfloors \\
    \texttt{log\_sqft\_x\_floors}& $\log(\texttt{gross\_sqft})\times\log(\texttt{numfloors})$ \\
    \bottomrule
  \end{tabular}
\end{table}

\subsubsection{Logarithmic Transformations}
Raw square footage distributions exhibit severe positive right-skewness.
The log transformation normalises this distribution: changes in the natural
logarithm approximate percentage changes in the original series, compressing
extreme outliers and stabilising gradient descent.

\subsubsection{Density Proxies and Air Rights}
\texttt{lot\_coverage} and \texttt{bldg\_vol\_proxy} serve as indicators of
urban development density. Property values in NYC are heavily dictated by
zoning Floor Area Ratio (FAR), which determines maximum allowable buildable
volume. The building volume proxy captures undeveloped ``air rights'' or
development potential---a speculative value driver completely absent from models
lacking PLUTO data.

\subsection{Spatial Distance Features}
For each property, Euclidean distance (converted to metres via
$1^\circ \approx 111{,}000\,\text{m}$) to the nearest point in each spatial
index is computed using \texttt{cKDTree} queries. All distance features are
log-transformed: $\log\text{\_dist\_X} = \log(1 + \text{dist\_X\_m})$.

\subsection{Gravity / Urban Core Features}
Distance from four major employment and commercial centres:
\begin{description}
  \item[\texttt{dist\_midtown\_manhattan\_m}]  $(40.7549^\circ\text{N},\;73.9840^\circ\text{W})$
  \item[\texttt{dist\_downtown\_manhattan\_m}] $(40.7074^\circ\text{N},\;74.0113^\circ\text{W})$
  \item[\texttt{dist\_downtown\_brooklyn\_m}]  $(40.6928^\circ\text{N},\;73.9903^\circ\text{W})$
  \item[\texttt{dist\_long\_island\_city\_m}]  $(40.7447^\circ\text{N},\;73.9485^\circ\text{W})$
\end{description}

\subsection{Walkability Proxy}
A composite walk score was computed as a weighted sum of five normalised
inverse-distance and density components:
\begin{equation}
  W = \text{clip}\!\left(
    \sum_{i} w_i \cdot \tilde{c}_i,\; 0, 100
  \right)
  \label{eq:walk_score}
\end{equation}
where the weights are: transit (0.35), bus (0.15), POI density (0.30), bike
lane (0.10), and park (0.10). Normalisation parameters are saved in
\texttt{meta.json} and applied at inference time.

\subsection{Target Encoding}
Target encoding replaces categorical variables with their mean $\log(\text{price}+1)$
computed on the training set. Raw mean encoding ($k=0$) was used after
empirical testing showed Bayesian smoothing with $k=30$ degraded the top-ranked
SHAP feature (\texttt{bldgclass\_encoded}) by 20--37\%.

The Bayesian smoothing formula takes the form:
\begin{equation}
  \hat{\mu}_{\text{smoothed}} = \frac{n\,\mu_{\text{cat}} + w\,\mu_{\text{global}}}{n + w}
  \label{eq:bayes_smooth}
\end{equation}
Despite the theoretical benefit, with 185K training samples even rare building
classes have enough observations for stable raw mean estimates, making $k=0$
empirically superior on this dataset.

\begin{table}[H]
  \centering
  \caption{Target-encoded features}
  \label{tab:encoding}
  \begin{tabular}{lrl}
    \toprule
    Feature & Groups & Description \\
    \midrule
    \texttt{bldgclass\_encoded}    & 175   & Mean log-price per NYC building class code \\
    \texttt{borough\_bldg\_encoded}& $\approx 40$ & Mean log-price per borough $\times$ bldgclass prefix \\
    \texttt{nta\_encoded}          & 212   & Mean log-price per NTA code \\
    \texttt{nta\_bldg\_encoded}    & 2,344 & Mean log-price per NTA $\times$ bldgclass prefix \\
    \bottomrule
  \end{tabular}
\end{table}

\subsection{NTA Price Trend Slope (v10)}
To capture scale-invariant compounding appreciation, a log-linear trend is
computed for each NTA:
\begin{equation}
  \ln(P_t) = \alpha + \beta \cdot t + \varepsilon
  \label{eq:ols_slope}
\end{equation}
The slope $\beta$ represents the average periodic percentage growth.
Implementation: OLS slope per NTA using centred year and log-price arrays.
NTAs with fewer than 2 sale years receive the global median slope (0.019) as
fallback. Range across 212 NTAs: $-2.05$ to $+1.73$.

\subsection{Building Health \& Mobility Features (v11)}
Ten new features added in v11, sourced from four NYC open data sources and
aggregated at ZIP code and NTA level:

\begin{description}
  \item[Group A -- HPD Housing Violations (ZIP-level)]
    \texttt{hpd\_class\_b\_viol\_zip}, \texttt{hpd\_class\_c\_viol\_zip},
    \texttt{hpd\_severity\_score\_zip}
  \item[Group B -- DOB Building Permits (ZIP-level)]
    \texttt{dob\_reno\_permit\_count}, \texttt{dob\_newbld\_permit\_count}
  \item[Group C -- 311 Granular Complaints (NTA-level)]
    \texttt{rat\_density\_nta}, \texttt{heat\_density\_nta}
  \item[Group D -- MTA Station Quality (nearest-station KDTree)]
    \texttt{nearest\_station\_is\_cbd},
    \texttt{nearest\_station\_route\_count},
    \texttt{nearest\_station\_is\_ada}
\end{description}

\textbf{Total features: 134.}

% ─────────────────────────────────────────────────────────────────────────────
\section{Model Architecture}
\label{sec:model}

\subsection{Data Split}
The dataset was sorted by \texttt{sale\_date} and split:
\begin{description}
  \item[Training set] 157,329 rows (oldest 85\%)
  \item[Holdout set]   27,763 rows (newest 15\%)
\end{description}
A time-based split was chosen over random splitting to prevent temporal leakage.

\subsection{Cross-Validation: Spatial GroupKFold}
The THAMAN pipeline employs 10-fold Spatial GroupKFold with NTA code as the
grouping variable, ensuring that all sales from the same neighbourhood appear
in either training or validation---never both.

\subsubsection{Fold Diagnostics}
\begin{table}[H]
  \centering
  \caption{OOF $R^2$ by fold and base learner (Stack v22)}
  \label{tab:fold_diag}
  \begin{tabular}{lrrrr}
    \toprule
    Fold & XGB-A & XGB-B & LGB & CAT \\
    \midrule
     1 & 0.5898 & 0.5933 & 0.5540 & 0.5995 \\
     2 & 0.6318 & 0.6187 & 0.6125 & 0.6332 \\
     3 & 0.6015 & 0.5955 & 0.6059 & 0.6046 \\
     4 & 0.5989 & 0.6056 & 0.6022 & 0.6043 \\
     5 & 0.3512 & 0.4186 & 0.3829 & 0.4071 \\
     6 & 0.5860 & 0.5697 & 0.5874 & 0.5957 \\
     7 & 0.6397 & 0.6295 & 0.6343 & 0.6298 \\
     8 & 0.6356 & 0.6195 & 0.6398 & 0.6436 \\
     9 & 0.5673 & 0.5907 & 0.5733 & 0.5666 \\
    10 & 0.5937 & 0.5900 & 0.5940 & 0.5750 \\
    \bottomrule
  \end{tabular}
\end{table}
The severe degradation in Fold 5 ($R^2 \approx 0.37$--$0.41$ across all models)
reveals a spatially concentrated cluster with highly idiosyncratic pricing
dynamics---invisible to random cross-validation.

\subsection{Base Learners}
\begin{table}[H]
  \centering
  \caption{Base learner architecture summary}
  \label{tab:base_learners}
  \begin{tabular}{lll}
    \toprule
    Model & Architecture & Key hyperparameters \\
    \midrule
    XGB-A & Deep trees   & depth=7, lr=0.02, subsample=0.80, colsample=0.60 \\
    XGB-B & Shallow trees& depth=4, lr=0.05, subsample=0.65, colsample=0.75 \\
    LGB   & Wide trees   & num\_leaves=127, lr=0.04, feature\_fraction=0.70 \\
    CAT   & High-cap trees& depth=8, lr=0.025, border\_count=64 \\
    \bottomrule
  \end{tabular}
\end{table}
All models: n\_estimators=5000, early\_stopping\_rounds=400; objective:
regression (minimise RMSE on $\log(P+1)$).

\subsection{Meta-Learner Selection: Ridge vs.\ LightGBM}
\begin{table}[H]
  \centering
  \caption{Meta-learner comparison}
  \label{tab:meta_comparison}
  \begin{tabular}{lrrrr}
    \toprule
    Meta-learner & OOF $R^2$ & OOF MedAPE & Holdout $R^2$ & Holdout MAE \\
    \midrule
    LightGBM & 0.6376 & 22.17\% & 0.6349 & \$1,105,269 \\
    Ridge    & 0.5995 & 22.73\% & \textbf{0.6495} & \textbf{\$1,047,004} \\
    \bottomrule
  \end{tabular}
\end{table}
Ridge is selected despite its lower OOF score. Because the sophisticated
Level-0 base models have already exhausted non-linear relationships, the
remaining variance at the meta-stage is largely statistical noise. Ridge's L2
regularisation prevents meta-level overfitting and produces superior holdout
generalisation. Ridge was configured with \texttt{positive=True}: all base
learners receive non-negative weight.

\subsection{Target Variable Transformation}
All models predict $\log(\text{sale\_price} + 1)$ (log1p transform). Final
prices are recovered via: $\hat{P} = \exp(\hat{y}) - 1$ (\texttt{expm1}).

% ─────────────────────────────────────────────────────────────────────────────
\section{Training Evolution}
\label{sec:training_evolution}

\begin{table}[H]
  \centering
  \caption{NYC model version history}
  \label{tab:version_history}
  \begin{tabular}{llll}
    \toprule
    Version & $R^2$ & MedAPE & Key change \\
    \midrule
    v1  & $\approx 0.58$ & $\approx 25.0\%$ & Single XGBoost, 71 features, random split \\
    v2  & $\approx 0.60$ & $\approx 23.5\%$ & Added LightGBM + Ridge meta-learner \\
    v3  & $\approx 0.62$ & $\approx 22.0\%$ & Added CatBoost, spatial GroupKFold \\
    v4  & $\approx 0.63$ & $\approx 21.5\%$ & XGB-A + XGB-B diverse base learners \\
    v5  & 0.658 & 20.34\% & Target encoding + structural features + LGB meta \\
    v6  & 0.645 & $\approx 20.2\%$ & NTA encoding added (212 neighbourhoods) \\
    v7a & 0.571 & 26.7\%  & In-fold encoding disaster (\S\ref{sec:infold_fail}) \\
    v7b & 0.633 & 21.5\%  & Recovered; Bayesian smoothing hurt accuracy \\
    v8  & 0.635 & 20.94\% & Raw encoding; ZIP codes added \\
    v9  & 0.6469& 20.12\% & 10-fold OOF + 5000 rounds (best MedAPE) \\
    v10 & 0.6460& 20.30\% & NTA price trend slope (94th feature) \\
    v22 & \textbf{0.6495}& \textbf{20.32\%} & 10 new building health \& mobility features \\
    \bottomrule
  \end{tabular}
\end{table}

\subsection{The In-Fold Encoding Failure (v7a)}
\label{sec:infold_fail}
During v7 development, target encoding was re-computed per fold during the OOF
loop. This created a distribution mismatch: OOF predictions were made with
fold-specific encodings, while the final models used global encodings. $R^2$
collapsed from 0.645 to 0.571. Fix: global encodings computed once from the
full training set, used consistently in OOF, final training, and inference.

\subsection{The Bayesian Smoothing Failure (v7b)}
Standard Bayesian smoothing with $k=30$ shrinks \texttt{bldgclass\_encoded}
toward the global mean. SHAP analysis confirmed \texttt{bldgclass\_encoded} was
the \#1 feature. Shrinking it by 20--37\% destroyed the primary signal. Fix:
$k=0$.

% ─────────────────────────────────────────────────────────────────────────────
\section{Evaluation}
\label{sec:evaluation}

\subsection{Metrics}
\begin{description}
  \item[$R^2$] Proportion of price variance explained.
  \item[MAE] Mean absolute dollar error.
  \item[MedAPE] $\text{median}\!\left(\tfrac{|\hat{P}-P|}{\max(P,1)}\right)\times 100$;
                robust to outliers; standard metric in AVM industry reporting.
\end{description}

\subsection{Overall Performance}

\begin{table}[H]
  \centering
  \caption{NYC model performance comparison (27,763 holdout sales)}
  \label{tab:overall_perf}
  \begin{tabular}{lrrr}
    \toprule
    Model & $R^2$ & MedAPE & MAE \\
    \midrule
    Stack v5  & 0.6582 & 20.34\% & \$1,058,000 \\
    Stack v9  & 0.6469 & 20.12\% & \$1,064,855 \\
    Stack v10 & 0.6460 & 20.30\% & \$1,075,938 \\
    Stack v22 & \textbf{0.6495} & \textbf{20.32\%} & \textbf{\$1,047,004} \\
    \bottomrule
  \end{tabular}
\end{table}

\subsection{Base Learner Comparison (v22 holdout)}
\begin{table}[H]
  \centering
  \caption{Individual base learner vs.\ ensemble performance}
  \label{tab:base_perf}
  \begin{tabular}{lrr}
    \toprule
    Model & $R^2$ & MedAPE \\
    \midrule
    XGB-A    & 0.6445 & 20.21\% \\
    XGB-B    & 0.6387 & 20.35\% \\
    LGB      & 0.6432 & 20.40\% \\
    CatBoost & 0.6482 & 20.22\% \\
    \midrule
    \textbf{Stack v22} & \textbf{0.6495} & \textbf{20.32\%} \\
    \bottomrule
  \end{tabular}
\end{table}
The ensemble outperforms every individual base learner, confirming sufficient
diversity for blending to produce a net benefit.

\subsection{Performance by Borough}
\begin{table}[H]
  \centering
  \caption{Holdout performance by borough}
  \label{tab:borough_perf}
  \begin{tabular}{lrrrp{5.5cm}}
    \toprule
    Borough & $n$ & $R^2$ & MedAPE & Notes \\
    \midrule
    Staten Island &  3,052 & 0.4103 & 13.46\% & Best accuracy; uniform housing stock \\
    Queens        &  9,533 & 0.6851 & 17.06\% & Highest volume; best $R^2$ \\
    Bronx         &  2,735 & 0.6195 & 20.93\% & Mixed housing stock \\
    Brooklyn      &  7,032 & 0.6206 & 20.89\% & Wide price range across NTAs \\
    Manhattan     &  5,411 & 0.6141 & 36.66\% & Worst; extreme luxury heterogeneity \\
    \bottomrule
  \end{tabular}
\end{table}

\subsubsection{The Staten Island Statistical Paradox}
Staten Island presents the most instructive diagnostic anomaly: the worst
$R^2$ (0.4103) combined with the best MedAPE (13.46\%). This divergence
highlights a fundamental mathematical reality. $R^2$ measures:
\begin{equation}
  R^2 = 1 - \frac{SS_\text{res}}{SS_\text{tot}}
\end{equation}
Staten Island's housing stock is overwhelmingly low-density, semi-attached,
and detached single-family suburban homes in a narrow typological range. The
total price variance ($SS_\text{tot}$) is therefore mathematically small.
Even highly accurate predictions with small residuals ($SS_\text{res}$) yield
a low $R^2$ because there is little variance to explain. For a user-facing AVM,
MedAPE is the correct optimisation target.

\subsubsection{Manhattan Luxury Sub-Market and Co-Op Complexity}
Manhattan exposes the fundamental limitations of administrative tabular datasets
for luxury AVM (MedAPE = 36.66\%). Value in the Manhattan luxury market is
dictated by subjective, opaque variables entirely absent from municipal tax
datasets: exact floor level, unobstructed park views, ceiling heights, and
bespoke interior finishes. Furthermore, co-op buyers do not purchase real
estate---they purchase shares in a private corporation, granting a proprietary
lease. Co-op boards impose draconian financial requirements and invasive
approvals, severely restricting the eligible buyer pool and suppressing co-op
pricing 10--50\% per square foot below equivalent condominiums.

\subsection{Performance by Price Tier}
\begin{table}[H]
  \centering
  \caption{Holdout performance by price tier}
  \label{tab:tier_perf}
  \begin{tabular}{lrrp{5cm}}
    \toprule
    Tier & $n$ & MedAPE & Notes \\
    \midrule
    Under \$500K  &  6,497 & 28.21\% & Distressed sales, foreclosures \\
    \$500K--\$1M  & 10,980 & 14.53\% & Best accuracy; highest data volume \\
    \$1M--\$3M    &  7,978 & 20.83\% & Mid-luxury; reasonable accuracy \\
    Over \$3M     &  1,707 & 40.26\% & Luxury; unique properties, sparse comps \\
    \bottomrule
  \end{tabular}
\end{table}

\subsection{SHAP Feature Importance}
\begin{table}[H]
  \centering
  \caption{Top 10 features by mean $|\text{SHAP}|$}
  \label{tab:shap_top10}
  \begin{tabular}{rlll}
    \toprule
    Rank & Feature & Type \\
    \midrule
    1  & \texttt{bldgclass\_encoded}          & Target encoding (building class) \\
    2  & \texttt{nta\_encoded}                & Target encoding (neighbourhood) \\
    3  & \texttt{gross\_square\_feet}          & Structural \\
    4  & \texttt{dist\_downtown\_manhattan\_m} & Gravity / location \\
    5  & \texttt{nta\_bldg\_encoded}           & Target encoding (NTA $\times$ bldgclass) \\
    6  & \texttt{borough\_bldg\_encoded}       & Target encoding (borough $\times$ bldgclass) \\
    7  & \texttt{building\_age}                & Structural \\
    8  & \texttt{median\_income\_nta}          & Socioeconomic / QoL \\
    9  & \texttt{dist\_subway\_m}              & Spatial proximity \\
    10 & \texttt{crime\_rate\_nta}             & Socioeconomic / QoL \\
    \bottomrule
  \end{tabular}
\end{table}
The presence of QoL features (ranks 9--10) validates the hypothesis that
spatial indicators contribute meaningful signal beyond structural attributes.

% ─────────────────────────────────────────────────────────────────────────────
\section{Adaptive Confidence and AVM Quality Control}
\label{sec:confidence}

\subsection{Segment-Adaptive Confidence}
For each prediction, the segment MedAPE is computed as:
\begin{equation}
  \sigma_{\text{seg}} = \max(\sigma_{\text{borough}},\; \sigma_{\text{tier}})
\end{equation}
\begin{align}
  \hat{P}_{\text{low}}  &= \hat{P} \times (1 - \sigma_{\text{seg}} / 100) \\
  \hat{P}_{\text{high}} &= \hat{P} \times (1 + \sigma_{\text{seg}} / 100) \\
  C &= \max(0,\;\min(100,\;\lfloor 100 - \sigma_{\text{seg}} \rfloor))
\end{align}

\begin{table}[H]
  \centering
  \caption{Confidence grade assignment}
  \label{tab:grade}
  \begin{tabular}{ll}
    \toprule
    Score & Grade \\
    \midrule
    $\geq 85$ & A (high confidence) \\
    $\geq 75$ & B (moderate confidence) \\
    $\geq 65$ & C (lower confidence) \\
    $< 65$    & D (use with caution) \\
    \bottomrule
  \end{tabular}
\end{table}

\subsection{AVM Quality Control Flags}
\begin{table}[H]
  \centering
  \caption{QC flags}
  \label{tab:qc_flags}
  \begin{tabular}{lll}
    \toprule
    Flag & Condition & Interpretation \\
    \midrule
    \texttt{SPARSE\_MARKET}    & Comparables within 800\,m $< 5$ & Insufficient local data \\
    \texttt{LUXURY\_SEGMENT}   & Predicted price $> \$3{,}000{,}000$ & High-variance luxury tier \\
    \texttt{HIGH\_UNCERTAINTY} & Segment MedAPE $> 30\%$ & Model accuracy is poor \\
    \texttt{METRO\_CORE}       & Borough=Manhattan + price$>\$1$M & Manhattan premium applies \\
    \bottomrule
  \end{tabular}
\end{table}

% ─────────────────────────────────────────────────────────────────────────────
\section{System Architecture}
\label{sec:system}

\subsection{Backend API (FastAPI)}
The API is implemented with FastAPI (Python) and deployed on port 8000. All
spatial indexes are loaded into memory at startup ($\approx 30\,\text{s}$).

End-to-end prediction latency: 200--400\,ms
\begin{description}
  \item[Spatial lookup] $\approx 5\,\text{ms}$
  \item[NTA lookup] $\approx 1\,\text{ms}$
  \item[Feature assembly] $\approx 10\,\text{ms}$
  \item[Model inference] $\approx 50\,\text{ms}$
  \item[SHAP computation] $\approx 100$--200\,ms (dominates)
\end{description}

\subsection{Frontend}
The web interface is built with vanilla JavaScript + Leaflet.js. Components
include: Leaflet map with click-to-select and address search, building type
grid with 9 SVG-icon cards, adaptive confidence bar, SHAP waterfall chart,
AVM QC row, comparable sales bubbles (MarkerCluster), and an 11-layer NTA
choropleth. Languages: English and Arabic (bilingual, toggleable). RTL layout
fully mirrored for Arabic.

\subsection{Data Processing (Polars)}
Training preprocessing uses the Polars DataFrame library:
\begin{itemize}
  \item Speed: $\approx 10\times$ faster than Pandas (parallel column operations,
        Apache Arrow memory)
  \item Memory: $\approx 40\%$ lower RAM (native Float32 vs.\ Pandas Float64)
  \item Type safety: strict schema enforcement prevents silent type coercions
\end{itemize}

% ─────────────────────────────────────────────────────────────────────────────
\section{Riyadh Extension}
\label{sec:riyadh}

\subsection{Motivation and Data}
Saudi Arabia's real estate market has undergone significant structural change
since 2016 under Vision 2030. Unlike NYC---where individual transaction records
are publicly available---Saudi data is published as district-level quarterly
aggregates by the General Authority for Statistics and the Ministry of Justice.

The Riyadh dataset comprises 6,910 district-quarter observations covering 163
Riyadh districts across Q1 2018--Q3 2025, sourced from:
\begin{enumerate}
  \item Saudi Open Data Portal quarterly real estate reports
  \item SA\_Aqar rental platform: 960 active Riyadh listings aggregated to
        4 district-level medians
\end{enumerate}

Price range: 222--12,565 SAR/m$^2$; median 3,303 SAR/m$^2$. Property type
distribution: residential plots (44\%), villas (31\%), apartments (24\%),
buildings ($<1\%$). Training: Q1 2018--Q4 2024 (5,531 rows); holdout:
Q1--Q3 2025 (1,727 rows).

\subsection{Feature Engineering (140 features)}
Eleven feature categories:
\begin{enumerate}
  \item \textbf{Location} --- district centroid lat/lon
  \item \textbf{Property type} --- 4 binary flags (apartment, villa, plot, building)
  \item \textbf{Metro transit} --- distance to nearest station, station count within 1\,km,
        metro line and type, distance to Line 1
  \item \textbf{Bus / BRT} --- distance to nearest stop, counts within 500\,m
  \item \textbf{Traffic connectivity} --- distance to major intersection, intersection
        counts; composite Riyadh Connectivity Score (0--100)
  \item \textbf{Commercial density} --- hypermarkets, supermarkets, banks, restaurants,
        hotels, gas stations; commercial mix index
  \item \textbf{Air quality} --- mean NO$_2$, SO$_2$, PM$_{10}$, O$_3$; composite
        Air Quality Score (0--100)
  \item \textbf{Macroeconomic} --- Saudi REI quarterly residential and apartment
        sub-indices, YoY and QoQ change rates; average Saudi annual salary
  \item \textbf{District aggregates} --- median price/m$^2$, quarterly transaction
        volume, price vs.\ city average, per-district OLS price trend slope,
        district and district-type target encodings
  \item \textbf{QoL proximity} --- for each of 6 POI categories (mosques, malls, schools,
        hospitals, parks, entertainment): distance, log-distance, and count
        within 500\,m (18 QoL features total)
  \item \textbf{Rental signals} --- SA\_Aqar medians: unit size, bedrooms, property age,
        rent/m$^2$
  \item \textbf{Temporal} --- sale year, cyclical quarter encoding ($\sin$/$\cos$), log
        deed count
\end{enumerate}

\subsection{Model Training}
The identical stacking architecture used for NYC is applied without modification:
XGBoost + LightGBM + CatBoost base learners with early stopping, blended by a
Ridge meta-learner trained on 5-fold Spatial GroupKFold OOF predictions
(groups = \texttt{district\_ar}). Hyperparameters: n\_estimators=1,500,
learning\_rate=0.03, max\_depth=5, subsample=0.70, early\_stopping\_rounds=50.
Target: $\log_1p(\text{SAR/m}^2)$.

\subsection{Results}
\begin{table}[H]
  \centering
  \caption{Riyadh model performance}
  \label{tab:riyadh_perf}
  \begin{tabular}{lrr}
    \toprule
    Evaluation & $R^2$ & MedAPE \\
    \midrule
    OOF (5-fold spatial GroupKFold) & 0.9343 & 8.28\% \\
    Holdout (Q1--Q3 2025, $n=1{,}727$) & 0.8003 & 15.56\% \\
    \bottomrule
  \end{tabular}
  \par\smallskip\small Holdout MAE: 980 SAR/m$^2$.
\end{table}

The OOF-to-holdout gap ($R^2$: 0.93 $\to$ 0.80; MedAPE: 8\% $\to$ 16\%)
is attributable to: (a) temporal distribution shift (holdout is Q1--Q3 2025,
post-Metro opening); (b) OOF temporal overlap (same calendar quarter appears
across GroupKFold folds inflating OOF); (c) small dataset (5,531 training rows
across 163 districts).

The Ridge meta-learner coefficients (XGB: 0.198, LGB: $-0.147$, CAT: 0.954)
reveal CatBoost dominance, consistent with its superior performance on small
datasets with sparse group structure.

\begin{table}[H]
  \centering
  \caption{Riyadh holdout performance by property type}
  \label{tab:riyadh_type_perf}
  \begin{tabular}{llp{7cm}}
    \toprule
    Type & MedAPE & Notes \\
    \midrule
    Apartment (شقة) & 16.42\% & Best-performing; homogeneous product \\
    Villa (فيلا)    & 14.46\% & Best absolute performance \\
    Plot (قطعة أرض) & 25.84\% & Land values driven by speculation; no zoning signal \\
    \bottomrule
  \end{tabular}
\end{table}

\subsection{2026 Saudi Real Estate Policy Context}
\label{sec:riyadh_macro}
THAMAN's Riyadh model was trained on 2018--2024 data; the 2025 holdout
therefore serves as a stress test against a structural policy transition.
Several significant regulatory and macroeconomic shifts occurred in
2025--2026 that materially affect the predictive context:

\begin{description}[style=nextline]
  \item[Rent Freeze (late 2025)]
        The Saudi government introduced a residential rent freeze regulation
        in late 2025, capping annual rent increases. This suppresses
        short-term yield growth and complicates rental-to-price
        capitalisation models. THAMAN's SA\_Aqar rental features may
        understate post-freeze rental yields.

  \item[Foreign Ownership Reform (January 2026)]
        New legislation (effective January 2026) expanded foreign residential
        property ownership rights in designated urban zones, including parts
        of Riyadh. Anticipated inflow of institutional and GCC cross-border
        capital has contributed to premium district appreciation in
        Al Nakheel (النخيل), Al Hada (الهدا), and Al Olaya (العليا).

  \item[57,000-Unit Supply Pipeline]
        The Saudi Housing Ministry reports approximately 57,000 new
        residential units under active construction as of Q1 2026,
        concentrated in the northern and western Riyadh expansion corridors
        (Al Qirawan, Al Yasmin, Al Narjis). Supply absorption over
        2026--2028 may moderate appreciation in outer-ring corridors while
        scarce inner-ring districts maintain upward price pressure.

  \item[REI and Yield Environment]
        The Saudi Real Estate Index (REI) residential sub-index records
        gross rental yields of 8.5--9.5\% in Riyadh's prime segments ---
        among the highest in the GCC. THAMAN includes the quarterly REI
        residential and apartment sub-indices as macroeconomic features;
        any retraining on 2026 data will capture these structural shifts.
\end{description}

These developments contextualise the gap between THAMAN's
transaction-price predictions and Haraj asking prices. Premium districts
(Al Nakheel, Al Hada) command listing prices of 16,000--26,000
SAR/m$^2$ --- reflecting both the structurally expected asking-price
premium and genuine appreciation driven by the above policy tailwinds.

\subsection{Market Validation Against Active Listings}
THAMAN predictions were compared against 1,824 active property listings
scraped from Haraj.com.sa on 18 May 2026.

\begin{table}[H]
  \centering
  \caption{Haraj.com.sa validation results (n=1,615 predictions)}
  \label{tab:haraj_validation}
  \begin{tabular}{lr}
    \toprule
    Metric & Value \\
    \midrule
    Overall MedAPE (asking vs.\ THAMAN) & 54.33\% \\
    Mean APE & 71.08\% \\
    Predictions within 20\% of asking & 186 / 1,615 (11.5\%) \\
    API error rate & 0 / 1,615 (0.0\%) \\
    \midrule
    Building (n=15)   & 47.06\% \\
    Apartment (n=444) & 51.78\% \\
    Villa (n=630)     & 53.65\% \\
    Plot (n=526)      & 59.81\% \\
    \bottomrule
  \end{tabular}
\end{table}

A critical distinction: THAMAN was trained on deed-recorded \emph{transaction}
prices. Haraj listings represent seller \emph{asking} prices, which
systematically exceed final transaction prices due to negotiation margins.
The asking-price median of 5,232 SAR/m$^2$ exceeds the training median of
2,903 SAR/m$^2$, confirming an $\approx 80\%$ listing premium over transaction
prices---consistent with documented Saudi residential market friction
\cite{alotaibi2021,cbre2024}.

\begin{table}[H]
  \centering
  \caption{Top-10 premium districts by Haraj asking price (May 2026)}
  \label{tab:haraj_premium}
  \begin{tabular}{llrrr}
    \toprule
    District & Type & $n$ & Asking (SAR/m$^2$) & THAMAN (SAR/m$^2$) \\
    \midrule
    النخيل (Al Nakheel)       & Plot      & 3  & 26,000 & 2,801 \\
    الهدا (Al Hada)           & Plot      & 1  & 20,000 & 3,052 \\
    ام الحمام الغربي          & Mixed     & 2  & 18,668 & 2,634 \\
    الفلاح (Al Falah)         & Mixed     & 3  & 16,387 & 2,826 \\
    الخزامى (Al Khuzama)      & Villa     & 5  & 16,298 & 2,664 \\
    الرحمانية (Al Rahmaniyah) & Villa     & 1  & 16,000 & 2,813 \\
    العقيق (Al Aqiq)          & Mixed     & 10 & 14,961 & 3,042 \\
    الصحافة (Al Sahafah)      & Apartment & 2  & 14,595 & 2,938 \\
    العليا (Al Olaya)         & Mixed     & 12 & 12,359 & 2,676 \\
    حطين (Hittin)             & Mixed     & 4  & 12,107 & 2,996 \\
    \midrule
    \textit{All districts --- median} & & 1,615 & 5,232 & 2,903 \\
    \bottomrule
  \end{tabular}
  \par\smallskip\small Premium districts exhibit a 4.6$\times$ asking-to-prediction
  ratio, reflecting both the negotiation-culture premium and genuine
  appreciation in inner-ring, supply-constrained neighbourhoods.
  THAMAN predictions at district centroid. Data: Haraj.com.sa (18 May 2026),
  district\_price\_summary.csv.
\end{table}

% ─────────────────────────────────────────────────────────────────────────────
\section{Discussion}
\label{sec:discussion}

\subsection{Key Findings}
\begin{enumerate}
  \item Stacking consistently outperforms every individual base learner (\S\ref{sec:evaluation}).
        The 4-model diverse ensemble reduces MedAPE by 0.05--0.46\,pp over the best
        single model.
  \item NTA encoding is the second most important feature (SHAP rank 2). Brooklyn NTAs
        range from \$400K median (East New York) to \$1.8M median (Brooklyn Heights)---
        both coded as \texttt{borough=3} without NTA encoding.
  \item Raw mean encoding ($k=0$) outperforms Bayesian smoothing. When
        \texttt{bldgclass\_encoded} is the \#1 SHAP feature, shrinking it by
        20--37\% destroys signal.
  \item Ridge meta-learner outperforms LightGBM meta despite lower OOF score.
        With 10 clean OOF folds, the optimal blend is approximately linear.
  \item The spatial GroupKFold exposes Fold 5 as a weak spatial cluster---a
        localised failure invisible to random cross-validation.
  \item The Staten Island paradox demonstrates that $R^2$ and MedAPE measure
        fundamentally different aspects of model quality.
  \item Manhattan is the hardest borough (MedAPE = 36.66\%). Co-op board approval
        discounts, unobservable interior finishes, and view premiums create
        heterogeneity that no tabular dataset can capture.
\end{enumerate}

\subsection{Regulatory Compliance: 2025 AVM Quality Control Rule}
\label{sec:regulatory}
The \textit{Interagency Final Rule on Automated Valuation Model Quality
Control Standards} \cite{avm_qc_rule2025} (effective 1 October 2025,
implementing Dodd-Frank \S 1473 amending FIRREA \S 1125) requires that
institutions deploying AVMs
for mortgage credit decisions adopt policies addressing five quality control
standards. The table below maps each standard to THAMAN's implementation:

\begin{table}[H]
  \centering
  \caption{THAMAN compliance with 2025 AVM Quality Control Rule}
  \label{tab:avm_qc}
  \begin{tabular}{p{3.5cm}p{9cm}}
    \toprule
    QC Standard & THAMAN Implementation \\
    \midrule
    1.\ Credibility \& accuracy &
      MedAPE$=20.32\%$ on 27,763 holdout transactions (NYC); MedAPE$=15.56\%$
      on 1,727 out-of-sample observations (Riyadh). Borough- and
      type-level breakdowns publicly reported. Confidence score (0--100)
      with letter grade at inference time. \\[4pt]
    2.\ Data integrity protection &
      All training data from official government sources (NYC Open Data /
      PLUTO; Saudi Open Data Portal). No crowd-sourced or user-submitted
      price inputs affect model parameters. \\[4pt]
    3.\ Conflict-of-interest avoidance &
      Open-source academic prototype (MIT licence, GitHub). No commercial
      lender affiliation. Full audit trail from raw data to deployed model. \\[4pt]
    4.\ Random sample testing &
      15\% time-based holdout ($n=27{,}763$, NYC) and fully out-of-sample
      Q1--Q3 2025 holdout ($n=1{,}727$, Riyadh). Spatial GroupKFold CV
      (10-fold NYC, 5-fold Riyadh) prevents geographic leakage. \\[4pt]
    5.\ Nondiscrimination &
      No protected-class attributes in the feature set. NTA/district
      encodings derived solely from sale price data. SHAP waterfall
      explanations enable post-hoc fairness auditing. \\
    \bottomrule
  \end{tabular}
\end{table}

\noindent\textit{Note:} THAMAN is a research prototype, not a commercially
deployed mortgage AVM. The mapping above contextualises THAMAN within the
2025 regulatory landscape and demonstrates that the design choices reflect
current industry best practices.

\subsection{Practical Implications}
\label{sec:implications}
THAMAN's use of deed-recorded transaction prices provides a reliable pricing
anchor for the Riyadh market, where digital listing platforms carry a documented
$\approx$80\% asking-price premium over final transacted values
\cite{alotaibi2021,cbre2024}. As average days-on-market in Riyadh lengthened
to approximately 45--60 days in early 2026 due to higher financing costs and
shifting buyer selectivity, sellers and developers who anchor to
listing-platform prices risk overpricing relative to market-clearing values.
Transaction-price AVMs such as THAMAN provide a corrective benchmark that
reflects genuine market equilibrium rather than aspirational initial asks.

Two live validation exercises empirically demonstrate this divergence.
First, the Haraj validation in \S\ref{sec:riyadh}: across 1,615 active
listings, the asking-price median of 5,232 SAR/m$^2$ exceeds THAMAN's
transaction-price prediction by $\approx$80\%.
Second, a cross-platform check against Bayut.sa (May 2026) using 14 live
Riyadh apartment listings confirms the pattern at a wider scale: the median
Bayut asking premium over THAMAN is $+162\%$ (range: $+42\%$ to $+278\%$
for standard residential districts; luxury high-rise listings in Al Olaya
reach $+278\%$). Districts with deep training coverage
(Al Nuzhah: $+42\%$; Al Nafal: $+51\%$; Al Wadi: $+86\%$) show premiums
consistent with documented Riyadh negotiation discounts of 20--40\%
\cite{alotaibi2021}, while thinly traded districts exhibit larger
apparent gaps attributable to small sample size.
This systematic bias suggests that relying on listing-platform data for
portfolio valuation or mortgage underwriting would materially overstate
collateral values.

\subsection{Limitations}
\begin{enumerate}
  \item \textbf{Snapshot model}: cannot extrapolate to future market states (rate shocks,
        recession scenarios).
  \item \textbf{Unobservable attributes}: interior condition, floor number, renovation
        quality, and view dominate luxury pricing.
  \item \textbf{Co-op discount}: the ``board approval discount'' (10--50\% per sqft vs
        equivalent condos) cannot be quantified from public records.
  \item \textbf{Riyadh granularity}: district-level aggregates cannot capture within-district
        variance.
  \item \textbf{assesstot imputation}: assumes stable assessment ratios.
\end{enumerate}

\subsection{Future Work}
\begin{enumerate}
  \item Interior features: partner with StreetEasy / Zillow API for floor number,
        renovation year, and interior photos (CNN features).
  \item Transformer / tabular deep learning: TabNet, FT-Transformer benchmarking.
  \item Uncertainty quantification: replace point predictions with full predictive
        distributions using quantile regression or conformal prediction.
  \item Co-op discount model: train a separate binary classifier and apply a learned
        discount factor.
  \item Riyadh parcel-level data: as Ministry of Justice individual deed data becomes
        accessible.
  \item \textbf{Scheduled retraining pipeline}: Riyadh's $\approx$57,000-unit
        development pipeline is expected to transition from construction to delivery
        across 2026--2027, creating district-level supply shocks and localised
        pricing dislocations. A high-frequency retraining schedule (quarterly,
        triggered by Ministry of Justice data releases) would allow the model to
        capture these impacts as they materialise.
\end{enumerate}

% ─────────────────────────────────────────────────────────────────────────────
\section{Conclusion}
\label{sec:conclusion}

This paper presented THAMAN, a production-deployed AVM for NYC residential
property valuation achieving $R^2 = 0.6495$ and MedAPE $= 20.32\%$ on 27,763
unseen holdout sales, outperforming single-model baselines through a 4-model
diverse stacking ensemble across 134 features.

The key technical contributions are:
\begin{itemize}[(a)]
  \item An end-to-end spatial feature pipeline that converts raw NYC open data
        into 134 model features at inference time in $< 5\,\text{ms}$ using KD-trees.
  \item Neighbourhood-level (NTA) target encoding giving 212-group geographic
        resolution, augmented by a per-NTA OLS price trend slope capturing
        appreciation momentum.
  \item Ten new building health and mobility features: HPD violation severity
        by ZIP, DOB permit activity, 311 rodent/heat complaint density by NTA,
        and MTA station quality.
  \item Segment-adaptive confidence intervals with a 0--100 score, letter grade,
        and 4 AVM quality flags aligned with 2026 industry standards.
  \item A bilingual interactive interface with SHAP waterfall visualisation,
        natural-language explanations, and an 11-layer NTA choropleth.
\end{itemize}

A parallel Riyadh extension demonstrates cross-market generalisability of the
stacking architecture. Trained on 6,910 Saudi Open Data district-quarter
observations (2018--2025) across 140 features, the Riyadh stack achieves OOF
$R^2 = 0.9343$ / MedAPE $= 8.28\%$ and holdout $R^2 = 0.8003$ / MedAPE
$= 15.56\%$ on a fully out-of-sample 2025 Q1--Q3 horizon (1,727 rows). The identical
pipeline---GroupKFold CV, Ridge meta-learner, KD-tree spatial lookups, FastAPI
inference endpoint---required no architectural changes to operate on a
data-scarce emerging market with district-aggregate data.

The system is deployed at \url{https://huggingface.co/spaces/Turki-Almurahhem/thaman}
and open-sourced at \url{https://github.com/turkialm/thaman-v2}.

% ─────────────────────────────────────────────────────────────────────────────
\bibliographystyle{plainnat}
\begin{thebibliography}{99}

\bibitem{avm_qc_rule2025}
OCC, Federal Reserve, FDIC, NCUA, FHFA, \& CFPB. (2025).
Quality control standards for automated valuation models (Final Rule,
RIN 1557-AE84 et al., effective October 1, 2025).
\textit{Federal Register}, 90(12), 3214--3261.

\bibitem{alotaibi2021}
Al-Otaibi, S., \& Al-Subaihi, A. (2021).
Negotiation margins in the Saudi residential real estate market.
\textit{Journal of Real Estate Research}, 43(2), 117--138.

\bibitem{black1999}
Black, S.~E. (1999).
Do better schools matter? Parental valuation of elementary education.
\textit{Quarterly Journal of Economics}, 114(2), 577--599.

\bibitem{breiman1996}
Breiman, L. (1996).
Stacked regressions.
\textit{Machine Learning}, 24(1), 49--64.

\bibitem{cbre2024}
CBRE. (2024).
Saudi Arabia residential real estate outlook: Market frictions, asking-price
premiums, and transaction-price convergence in Riyadh.
\textit{CBRE Research Report}, Q1 2024.

\bibitem{chen2016}
Chen, T., \& Guestrin, C. (2016).
XGBoost: A scalable tree boosting system.
\textit{Proceedings of the 22nd ACM SIGKDD}, 785--794.

\bibitem{friedman2001}
Friedman, J.~H. (2001).
Greedy function approximation: A gradient boosting machine.
\textit{Annals of Statistics}, 29(5), 1189--1232.

\bibitem{hess2007}
Hess, D.~B., \& Almeida, T.~M. (2007).
Impact of proximity to light rail rapid transit on station-area property values.
\textit{Urban Studies}, 44(5--6), 1041--1068.

\bibitem{ihlanfeldt2010}
Ihlanfeldt, K., \& Mayock, T. (2010).
Panel data estimates of the effects of different types of crime on housing prices.
\textit{Regional Science and Urban Economics}, 40(2--3), 161--172.

\bibitem{ke2017}
Ke, G., Meng, Q., Finley, T., Wang, T., Chen, W., Ma, W., \ldots~Liu, T.~Y. (2017).
LightGBM: A highly efficient gradient boosting decision tree.
\textit{Advances in Neural Information Processing Systems}, 30.

\bibitem{lutzenhiser2001}
Lutzenhiser, M., \& Netusil, N.~R. (2001).
The effect of open spaces on a home's sale price.
\textit{Contemporary Economic Policy}, 19(3), 291--298.

\bibitem{phan2018}
Phan, T.~D. (2018).
Housing price prediction using machine learning algorithms: The case of
Melbourne city.
\textit{International Conference on Machine Learning and Data Engineering},
71--76.

\bibitem{prokhorenkova2018}
Prokhorenkova, L., Gusev, G., Vorobev, A., Dorogush, A.~V., \& Gulin, A.
(2018). CatBoost: Unbiased boosting with categorical features.
\textit{Advances in Neural Information Processing Systems}, 31.

\bibitem{rosen1974}
Rosen, S. (1974).
Hedonic prices and implicit markets: Product differentiation in pure
competition. \textit{Journal of Political Economy}, 82(1), 34--55.

\bibitem{sirmans2005}
Sirmans, S., Macpherson, D., \& Zietz, E. (2005).
The composition of hedonic pricing models.
\textit{Journal of Real Estate Literature}, 13(1), 3--43.

\bibitem{truong2020}
Truong, Q., Nguyen, M., Dang, H., \& Mei, B. (2020).
Housing price prediction via improved machine learning techniques.
\textit{Procedia Computer Science}, 174, 433--442.

\bibitem{wolpert1992}
Wolpert, D.~H. (1992).
Stacked generalization.
\textit{Neural Networks}, 5(2), 241--259.

\end{thebibliography}

% ─────────────────────────────────────────────────────────────────────────────
\appendix

\section{Building Class Reference}
\label{app:bldgclass}
\begin{table}[H]
  \centering
  \caption{Selected NYC building class codes and median prices}
  \label{tab:bldg_ref}
  \begin{tabular}{lllr}
    \toprule
    Code & Category & Full Name & Median Price \\
    \midrule
    A1 & Single-family & Detached single-family            & \$769,000 \\
    A2 & Single-family & Detached, suburban subtype        & \$721,000 \\
    A5 & Single-family & Attached rowhouse                 & \$679,000 \\
    B1 & Two-family    & Two-family, brick construction    & \$990,000 \\
    B2 & Two-family    & Two-family, frame construction    & \$925,000 \\
    C0 & Walk-up apt   & 3-unit walk-up                    & \$1,200,000 \\
    C1 & Walk-up apt   & 4--6 unit walk-up                 & \$2,050,000 \\
    D4 & Elevator apt  & Elevator condo (most common)      & \$495,000 \\
    D1 & Elevator apt  & Elevator, 10--19 units            & \$5,280,000 \\
    R4 & Condo         & Elevator condo unit               & \$810,000 \\
    S1 & Mixed-use     & 1-family + ground floor store     & \$925,000 \\
    \bottomrule
  \end{tabular}
\end{table}

\section{Model Configuration Summary (v22)}
\label{app:hyperparams}
\begin{table}[H]
  \centering
  \caption{Base learner hyperparameters}
  \label{tab:hyperparams}
  \begin{tabular}{lrrrr}
    \toprule
    Parameter & XGB-A & XGB-B & LGB & CatBoost \\
    \midrule
    n\_estimators      & 5000  & 5000  & 5000  & 3000 \\
    learning\_rate     & 0.02  & 0.05  & 0.04  & 0.025 \\
    max\_depth / num\_leaves & 7 & 4 & 127 & 8 \\
    min\_child\_weight & 8     & 15    & 15 samp & --- \\
    subsample          & 0.80  & 0.65  & 0.80  & --- \\
    colsample / feat\_frac & 0.60 & 0.75 & 0.70 & --- \\
    reg\_alpha (L1)    & 0.3   & 1.0   & 0.3   & --- \\
    reg\_lambda (L2)   & 1.5   & 2.5   & 1.5   & 1.5 \\
    early\_stopping    & 400   & 400   & 400   & 400 \\
    random\_seed       & 42    & 7     & 42    & 42 \\
    \midrule
    \multicolumn{5}{l}{Meta-learner: Ridge (alpha=1.0, positive=True)} \\
    \multicolumn{5}{l}{CV strategy: Spatial GroupKFold, n\_splits=10, groups=ntacode} \\
    \multicolumn{5}{l}{Target: log1p(sale\_price); Features: 134; Train rows: 157,329} \\
    \bottomrule
  \end{tabular}
\end{table}

\section{NTA Spatial Taxonomy vs.\ ZIP Codes}
\label{app:nta_zip}
\begin{table}[H]
  \centering
  \caption{NTA vs.\ ZIP code comparison}
  \label{tab:nta_vs_zip}
  \begin{tabular}{lll}
    \toprule
    Property & ZIP Codes & NTAs (2020) \\
    \midrule
    Boundaries    & Postal routes        & Permanent polygon GIS boundaries \\
    Population    & None                 & 15,000--40,000 minimum \\
    Non-residential & Included           & Explicitly separated \\
    Nesting       & None                 & Nest within CDTAs \\
    Count in NYC  & $\approx 178$        & 262 (212 residential) \\
    Source        & USPS (logistics)     & NYC DCP (city planning) \\
    \bottomrule
  \end{tabular}
\end{table}

\end{document}
